\section{Syntax with a computable account of binding}\label{sec:binding}

Fix a finite signature, given by two lists of arities for functions and
predicates. Constants have function arity zero. The arguments also apply
schema by schema to a fixed signature given by primitive-recursive (PR)
symbol tests and arity functions. An expression at arity $n$ may refer to
the $n$ individuals $v_0,\ldots,v_{n-1}$. Terms are variables and
applications $f(\vec t)$ of the stipulated arity. Formulas are predicates
$R(\vec t)$, identities $s=t$, $\bot$, $A\wedge B$, $A\vee B$, $A\to B$,
$\forall B$ and $\exists B$. Under either quantifier the body has arity
$n+1$, with its newly bound slot at index zero. Write
$\neg A=A\to\bot$ and $A\leftrightarrow B=(A\to B)\wedge(B\to A)$.
An expression remains well scoped if an unused slot is appended at
the high end. A sentence has no free occurrences.

\begin{theorem}[Parsing, scope and alpha-conversion]\label{thm:alpha}
Well-formedness, the parse tree, the free-variable set, the binding
occurrence of each bound variable and alpha-equivalence of named
expressions all have PR algorithms. Parsing succeeds with a unique tree.
\end{theorem}
\begin{proof}
For a concrete serialization use prefix tokens carrying the constructor
and any index or symbol. Its arity fixes the number and sorts of its
children. Read the tokens backwards, pushing leaves and at each internal
node popping the prescribed number of children, checking sorts in their
specified order and pushing the resulting tree. Reject underflow,
a bad symbol or any final stack other than one element of the requested
sort. At each suffix the stack is its left-to-right list of complete
subtrees whenever that suffix can arise in a parse. Reading a parent is
then forced. Induction on the number of scanned tokens proves uniqueness
and proves that serializing any tree is inverted by this scan. To check
indices, traverse the resulting tree with the current binder depth $d$.
An index $j<d$ picks the binder $j$ positions back from the innermost.
Otherwise it denotes free slot $j-d$ and must satisfy $j-d<n$.
Mark those slots in an $n$-bit mask. This computes the exact support,
so closedness is decidable even when the input was supplied at positive
arity.

For named input keep a stack of active binder names, nearest first.
Look up each occurrence on that stack by bounded search. A hit gives its
position, and a miss gives the free name's index in a fixed table shared
by both expressions being compared. A common table is obtained by sorting
and removing duplicates from the free names in the two inputs. Renaming a
binder to a fresh name and changing exactly its bound occurrences preserves
these lookups. Conversely, assign to each binder in preorder the name
$M+1+j$ for its preorder number $j$, where $M$ exceeds all names in both
inputs. Both expressions can be renamed to this representative one binder
at a time without capture. Equal indexed trees now print identically.
The scans, stacks and bounded searches use at most the total token count
in steps per traversal or lookup. The fresh names are bounded by $M+1$
plus the larger token count.
\end{proof}

A substitution $\sigma:n\to m$ is a list of $n$ terms at arity $m$.
Let $E[\sigma]$ replace each free slot by its image. On terms the
clauses are $v_i[\sigma]=\sigma_i$ and
$f(\vec t)[\sigma]=f(\vec t[\sigma])$. For atomic formulas and
connectives apply the operation to the children. Put
$\uparrow t=t[v_1,\ldots,v_m]$ for a term at arity $m$, and use
the same simultaneous shift on formulas. At a quantifier use the lift
\[
 \sigma^+=(v_0,\uparrow\sigma_0,\ldots,\uparrow\sigma_{n-1}),
 \qquad (QB)[\sigma]=Q(B[\sigma^+]).
\]
Equivalently, at depth $d$ copy $v_j$ when $j<d$, and replace $v_{d+i}$
by $\sigma_i$ with every index increased by $d$. This last algorithm
defines shift on formulas directly. Write
$B\langle t\rangle=B[t,v_0,\ldots,v_{n-1}]$ for opening an
$(n+1)$-ary body at arity $n$.

\begin{theorem}[Simultaneous substitution]\label{thm:subst}
Substitution preserves scope, acts identically for the variable list
$\mathrm{id}_n$, and obeys
\[
  (E[\sigma])[\tau]=E[\sigma;\tau],\qquad
  (\sigma;\tau)_i=\sigma_i[\tau].
\]
Its exact free support is
$\operatorname{FV}(E[\sigma])
 =\bigcup_{i\in\operatorname{FV}(E)}\operatorname{FV}(\sigma_i)$.
Substitutions agreeing on the free support give identical results.
In particular
$(B\langle t\rangle)[\sigma]=B[\sigma^+]\langle t[\sigma]\rangle$.
\end{theorem}
\begin{proof}
At depth $d$ the copied bound indices are below $d$ and every inserted
index lies in $[d,d+m)$. A former free occurrence contributes exactly
the free variables of its replacement. Taking the union over the
occurrences gives the support formula, and comparison of the actual
insertions proves agreement. Replacing by the identity copies every
index. For composition consider a free occurrence $v_{d+i}$. After
the first operation its replacement is $\sigma_i$ shifted by $d$.
Every variable in this inserted term is therefore free relative to
the surrounding $d$ binders. Applying the second operation replaces
each original variable $v_j$ of $\sigma_i$ by $\tau_j$ shifted
by $d$. The result is precisely $\sigma_i[\tau]$ shifted by $d$.
Original bound occurrences and constructor nodes are copied on both
sides. This proves the equality of output trees, including when
$\sigma_i$ has no variables. For opening, the two composites send
$v_0$ to $t[\sigma]$ and $v_{i+1}$ to $\sigma_i$. Apply composition.

For explicit PR bounds, if $N$ counts input nodes and $L$ is the
largest image size (take $L=1$ for an empty list), the output has at
most $N\max(1,L)$ nodes. The depth is at most $N$, and a shifted
variable index is less than $m+N$ whenever one is inserted. A traversal
with a stack visits each original node and copies each selected image
within these bounds. The preceding arguments can be verified by a
bounded conjunction over occurrence addresses: trace each output block
to its original node, then use the computed input path and depth at that
node. This avoids induction with a universally quantified substitution
as induction formula.
\end{proof}

\section{Effective deduction and admissible proof operations}\label{sec:deduction}

Write $\Gamma\vdash_n A$ for a certificate in arity $n$ with an
ordered finite list $\Gamma$ of assumptions. Write $\Gamma,A$ for
appending an assumption, and $\uparrow\Gamma$ for shifting every
formula. A node records its sequent, a rule, its ordered child proofs
and the term or body used by its rule. All formulas and terms are
scope-checked. The two calculi below share the following rules.
\begin{itemize}
\item A leaf cites an assumption by position. Ex falso concludes
any $A$ from $\bot$. Conjunction is formed from both proofs and projected
to either component. Disjunction is introduced from either disjunct.
To eliminate $A\vee B$, prove $C$ under $\Gamma,A$ and under
$\Gamma,B$. Implication introduction discharges the last assumption
in $\Gamma,A\vdash B$ and elimination is modus ponens in a common
context.
\item Universal introduction has premise
$\uparrow\Gamma\vdash_{n+1}B$ and conclusion
$\Gamma\vdash_n\forall B$. Its elimination gives $B\langle t\rangle$
from $\forall B$. Existential introduction reverses this last pattern
with $\exists$ in the conclusion. Existential elimination has premises
$\Gamma\vdash_n\exists B$ and
$\uparrow\Gamma,B\vdash_{n+1}\uparrow C$, and conclusion
$\Gamma\vdash_n C$.
\item Reflexivity proves $t=t$. Equality elimination takes proofs
of $s=t$ and $B\langle s\rangle$ to a proof of $B\langle t\rangle$.
The unary formula context $B$ can also use the old $n$ slots.
\item A fresh, unused parameter may be dropped:
from $\uparrow\Gamma\vdash_{n+1}\uparrow A$ infer
$\Gamma\vdash_n A$. This implements the ordinary convention of
nonempty domains even in a signature without a closed term. For example,
in one parameter, reflexivity followed by existential introduction proves
the shift of $\exists x\,(x=x)$. The rule closes this proof.
\end{itemize}
Call this calculus $I$. Its extension $C$ adds one rule, from
$\neg\neg A$ infer $A$. The tags $I,C$ here describe object proofs.
All the algorithms and arguments about their codes use iPRA.

\begin{theorem}[Proof recognition and finite search]\label{thm:proofcheck}
For either calculus, the relation $\operatorname{Check}(p,n,\Gamma,A)$
is PR. The increasing bounded search
$S(b)=\bigvee_{p\le b}\operatorname{Check}(p,n,\Gamma,A)$ is PR,
and a positive result supplies a least certificate $p\le b$.
Consequently theoremhood in either presented calculus is recursively
enumerable.
\end{theorem}
\begin{proof}
Parse a proof code as a tree. Reject malformed input. At each node compare
its child sequents to the required instances, computing shifts, openings
and positional lookups. Mark it accepted exactly when these tests and
all child marks succeed. A bottom-up scan proves by induction over the
scan position that the accepted nodes are exactly the locally justified
subtrees with accepted children. The root must also match the four inputs.
Operations use a traversal bounded by input size and the PR substitution
algorithm. Bounded disjunction and least successful index are finite
folds. Enumerate proof codes by increasing number and output their
accepted sequents. Each actual proof occurs at its own code number.
\end{proof}

\begin{theorem}[Instance, context substitution, cut and deduction]\label{thm:graft}
Uniform PR operations on certificates implement all the following.
\begin{enumerate}
\item From $\Gamma\vdash_n A$ and $\sigma:n\to m$, obtain
$\Gamma[\sigma]\vdash_m A[\sigma]$.
\item Given $p:\Gamma\vdash_n A$ and, for every position $i$ of
$\Gamma$, a proof $q_i:\Delta\vdash_n\Gamma_i$, obtain
$\Delta\vdash_n A$. In particular weakening, exchange, contraction
and cut are admissible.
\item Certificates for $\Gamma,A\vdash B$ and for
$\Gamma\vdash A\to B$ can be converted in both directions.
A derivation uses only a computable sublist of its assumptions.
\end{enumerate}
\end{theorem}
\begin{proof}
For the first operation relabel every sequent and the rule's term and
formula annotations by $\sigma$. At an individual-variable binder use
$\sigma^+$ in the premise at larger arity. In initial contexts and
shifted conclusions we require
$(\uparrow E)[\sigma^+]=\uparrow(E[\sigma])$. The two composite
variable maps coincide by Theorem~\ref{thm:subst}. In instantiation
and in equality elimination we require that opening commute with
substitution, the last identity of that theorem. Thus the same rule
tag remains legal everywhere, including at parameter erasure.
Propositional binders merely retain their added assumption.

For the second operation, first implement change of context by a map
from each assumption position to a position carrying the same formula.
Change hypothesis leaves by the map. Under a discharged assumption
extend the map by matching the new last position to the new last
position. Under an individual binder the old position map also maps the
shifted formulas. For existential elimination extend that shifted map
by the new last position. All other steps copy their original rule.
This gives weakening and all positional structural rules directly.

Now graft $q_i$ at each original assumption leaf. On entering a
propositional discharge with new assumption $D$, weaken every current
$q_i$ to $\Delta,D$ and append the assumption leaf for $D$ to the
substitution list. At universal introduction or parameter erasure,
shift the current proofs by part 1. At the second child of existential
elimination do both operations, first shift and then append the proof
of its new local assumption. All premise sequents are thereby relabelled
as required. Cut is the case of one nontrivial graft and identity
assumption proofs for the other positions. Discharge gives the forward
deduction conversion in a single step. For the reverse direction weaken
the proof of $A\to B$ to $\Gamma,A$, append an assumption leaf for
$A$ and use modus ponens.

To extract the finite support, trace the cited hypothesis of each
assumption leaf back through the position maps to the root or to its
local discharge. Mark precisely the cited root positions. Delete the
other root positions throughout and reindex surviving citations and local
discharges. This yields a proof under the marked sublist. In particular,
a proof from a possibly infinite axiom collection, presented with a
membership certificate for each use, exhibits its own finite subcollection.
For a PR checker in that formulation the supplied membership certificates
must themselves have a PR verification relation.

Each algorithm visits the nodes of the original $p$ once, carrying a
list of current maps or proofs downward. A current graft has the original
rule-tree size of its $q_i$ or is a one-node local assumption: weakening
and shifting relabel vertices. The grafted tree thus has at most
$|p|(1+\max_i |q_i|)$ rule nodes, taking the empty maximum to be zero.
During relabelling the number of nested extensions is at most $|p|$.
Contexts are finite lists, with at most $|p|$ new entries on a path.
Repeated applications of a fixed PR bound for a single relabelling
are bounded by its primitive iteration $|p|$ times. More formally one
can run a forward pass producing the local arguments at all addresses
and a reverse pass producing the result. The assertion that every
processed address has passed its local equality tests is a single
bounded Boolean fold. This supplies quantifier-free induction for the
certificate transformations.
\end{proof}

\begin{proposition}[A Hilbert deduction transformer]\label{prop:hilbert-ded}
Consider a named-variable Hilbert presentation with modus ponens,
generalization on a variable absent from the free variables of all
undischarged assumptions, and axiom instances
\[
 \begin{gathered}
  K: U\to(V\to U),\qquad
  S:(U\to(V\to W))\to((U\to V)\to(U\to W)),\\
  (\forall x(U\to V))\to(U\to\forall x V)
       \qquad(x\notin\operatorname{FV}(U)).
 \end{gathered}
\]
Other fixed logical axioms may be included. From an $N$-line derivation
of $B$ from $\Gamma,A$, a PR algorithm produces a derivation of $A\to B$
from $\Gamma$ in at most $5N$ lines.
\end{proposition}
\begin{proof}
For each old line $D$ store the location of its new last line $A\to D$.
For an axiom or an assumption in $\Gamma$, copy $D$, append
$D\to(A\to D)$ and apply modus ponens. For the hypothesis $A$ use the
two $K$ instances $A\to((A\to A)\to A)$ and $A\to(A\to A)$.
Apply the matching $S$ instance and modus ponens twice to obtain
$A\to A$ in five lines. For modus ponens on $E$ and $E\to D$ use $S$
and two applications of modus ponens to their stored versions. If $D$
is $\forall x E$ by generalization, $x$ is absent from $A$ and from
$\Gamma$. Generalize the stored $A\to E$, append the displayed
quantifier axiom and apply modus ponens. All new references point
backwards. Checking the stored-location invariant for the bounded
list of earlier lines proves correctness by induction on the scan.
\end{proof}

\subsection{Equality and eliminating explicit definitions}

\begin{proposition}[Derived equality calculus]\label{prop:derived-eq}
In $I$ there are uniform proofs of symmetry, transitivity, congruence for
each function and predicate, and replacement of equals in every formula.
Conversely reflexivity, symmetry, transitivity and congruence for
the function and predicate symbols suffice for formula replacement.
\end{proposition}
\begin{proof}
Given $s=t$, substitute into $z=s$, starting with $s=s$, to get $t=s$.
From $s=t$ and $t=u$ substitute in $s=z$ to get $s=u$. Insert a single
varying argument into the context
$f(\ldots,s,\ldots)=f(\ldots,z,\ldots)$ and start with reflexivity.
For a predicate the context is $R(\ldots,z,\ldots)$. Change
successive arguments one at a time and compose the implications or
identities. Replacement for an arbitrary context is precisely the
elimination rule, and symmetry gives the reverse implication.

For the converse fix $s=t$ and a term context $u(z)$. A variable occurrence
of the hole uses $s=t$. Another variable uses reflexivity. At a function
node apply its congruence to the already obtained identities of its
arguments. For formula contexts construct both directions together.
At a relation use term congruence and predicate congruence in each
direction. At $u=v$ use transitivity and the two term identities,
reversing one by symmetry. Both occurrences of $\bot$ are identical.
At $\wedge$ convert each component. At $\vee$ reason by cases. At an
implication, first convert its antecedent backwards, apply the given
implication and convert the result forwards. The reverse proof exchanges
these two roles. Under $\forall$, open at a fresh parameter, apply the
body conversion with the shifted $s=t$, and generalize again. Under
$\exists$ open a witness in a local subproof, convert its body and
introduce the existential with that very witness. The conclusion omits
the local parameter, as required by elimination. The procedure recurses
over the finitely many input nodes. Store both conversions per node
and compose them by the graft operation. An explicit PR size bound is
obtained by iterating a bound for this construction at one node the
number of nodes in the input. Here a step combines finitely many
previously constructed proofs, at most the maximum symbol arity plus
a fixed number of rule templates, and uses the PR operations of
Theorem~\ref{thm:graft}.
\end{proof}

\begin{theorem}[Conservative predicate abbreviations]\label{thm:def-abbrev}
Let $D$ be an old-language formula in $k$ slots. Extend the language
by one fresh $k$-ary predicate $P$ and the universally closed axiom
$\forall\vec x\,(P(\vec x)\leftrightarrow D(\vec x))$ with the indicated
argument order. Every extended proof of an old sequent from old axioms
and this axiom yields an old proof of the sequent by a PR transformation,
in $I$ or in $C$. A finite ordered list of definitions referring only to
earlier ones can likewise be eliminated.
\end{theorem}
\begin{proof}
Expand $P(\vec t)$ into $D[\vec t]$ at every occurrence in sequents
and formula annotations, and recurse through the old constructors.
Expansion commutes with substitution: at $P(\vec t)$ this is the
composition identity of Theorem~\ref{thm:subst}. Under a binder
the same identity uses the lifted substitution. Expansion is the
identity on the old language. It therefore preserves every rule
test, including opening in equality elimination. Replace the expanded
defining axiom by its old proof: derive $D\to D$ by discharge twice
for the two conjuncts of $D\leftrightarrow D$, combine them and
universally introduce each parameter. In an arbitrary ambient context
one may first prove this closed theorem, then weaken and shift as needed.
Use grafting to replace all uses of the axiom. For a list, process the
definitions in reverse order. Iteration over the input list of the
uniform single-definition algorithm is PR.
\end{proof}

\section{Soundness, quantifier laws and small separating models}
\label{sec:deduction-soundness}

Here is an interpretation that checks which logical principles the two
certificate types support. A finite Kripke frame consists of a nonempty
finite list of worlds, a reflexive transitive order $\leq$, and nonempty
finite domains $D_w$ of natural-number labels, with $D_w\subseteq D_v$ for
$w\leq v$. At each world there are total function tables and relation
tables on its domain. Require that functions preserve their values on
earlier tuples and positive relation entries persist upwards. Equality is
equality of the individual labels. All these conditions are decidable
by bounded scans. For an assignment $a\in D_w^n$ interpret terms by table
lookup. Denote its value by $[t]_{w,a}$. Atomic forcing uses the given
tables, $\bot$ has value false, and $\wedge,\vee$ are the Boolean
conjunction and disjunction at the current world. The remaining clauses
are
\[
\begin{aligned}
 w,a\Vdash U\to V
    &\iff (\forall v\ge w)\,(v,a\Vdash U\ \Rightarrow\ v,a\Vdash V),\\
 w,a\Vdash\forall B
    &\iff (\forall v\ge w)(\forall e\in D_v)\,v,(e,a)\Vdash B,\\
 w,a\Vdash\exists B
    &\iff (\exists e\in D_w)\,w,(e,a)\Vdash B .
\end{aligned}
\]
The displayed quantifiers range over finite tables. All three tests
return bits. With one world these are precisely ordinary two-valued
semantics, also called Tarskian semantics here.

\begin{lemma}[Persistence and semantic substitution]\label{lem:krip-sub}
Term values are unchanged at later worlds on old assignments, and positive
forcing persists. For a substitution $\sigma:n\to m$ and an assignment
$a$ in $D_w^m$, put $b_i=[\sigma_i]_{w,a}$. Then
\[
  [t[\sigma]]_{w,a}=[t]_{w,b},\qquad
  w,a\Vdash A[\sigma]\ \Longleftrightarrow\ w,b\Vdash A.
\]
Truth depends only on the free support. In particular a sentence has the
same truth value under every assignment at a fixed world.
\end{lemma}
\begin{proof}
Variable lookup is unchanged at a larger world. The condition on
function tables propagates this assertion through each term node. For
forcing, persistence at atomic nodes follows by the condition on relations
and stability of term values. The cases $\bot,\wedge,\vee$ follow from their
definitions. For an implication or a universal formula, every extension
of a later world is already an extension of the earlier one. For an
existential formula its witness still belongs to the larger domain
and its body persists by the child assertion.

For term substitution, variable lookup gives the claimed equality on
variables and applying the same function entry gives it on a compound
term. Atom evaluation then agrees. Conjunction, disjunction and falsity
preserve agreement. At implication, in each $v\ge w$ the values of
$\sigma$ at $a$ are still the original $b$, so apply agreement to the
two subformulas there. At a quantifier, the image of $\sigma^+$ evaluated
at $(e,a)$ is $(e,b)$: its first component is the fresh slot and the
rest ignore that slot. This works at each $v\ge w$ for the universal
and at $w$ for the existential. Finally assignments agreeing on free
support have equal term values at all corresponding occurrences.
At depth $d$ extend both by the same $d$ chosen individual labels
along a common sequence of worlds. Each variable lookup in the
corresponding bodies then agrees, and the forcing clauses give the
same output. This also proves the sentence assertion.
\end{proof}

\begin{theorem}[Finite soundness]\label{thm:krip-check}
Every accepted $I$ proof preserves forcing at every world of every
coded finite Kripke model. Every accepted $C$ proof preserves truth in
every coded nonempty finite two-valued structure. Both claims are
uniform theorems of iPRA.
\end{theorem}
\begin{proof}
For each node, tabulate for every world and every assignment the assertion
that if all its assumptions are forced then its conclusion is forced.
An assumption is a lookup. A false conclusion from an assumed proof of
$\bot$ poses no case, because a forced $\bot$ is the false bit. The
conjunction and disjunction rules read the corresponding bits. For
disjunction elimination choose the true disjunct and use its branch with
that assumption appended. For implication introduction pass to any $v\ge w$
forcing the antecedent. Persistence takes every old assumption to $v$,
so the branch gives the consequent. For elimination use the implication
test at $v=w$ and the true antecedent.

For universal introduction take arbitrary $v\ge w$ and $e\in D_v$.
Each old assumption, shifted, is true at $(v,(e,a))$ by persistence and
the semantic substitution lemma. The premise gives the body there.
For elimination choose $v=w$ and $e=[t]_{w,a}$, then use semantic
substitution. An existential introduction uses the same $e$ as witness.
For elimination scan for the first witness $e$ of the forced major premise
at $w$. At $(w,(e,a))$ the shifted old assumptions and the body all hold,
so the minor premise gives the shifted conclusion. Ignoring the fresh
slot gives the desired conclusion at $(w,a)$. For parameter erasure
use the least label of nonempty $D_w$ and the same shift identity.
Reflexivity uses identity of labels. In equality elimination, $s,t$
have the same value $e$, so both opened formulas have the truth of the
same body at $(e,a)$. In a one-world model implication is the usual
Boolean implication. Its double-negation operation is the identity
on each of the two bits, verifying the additional $C$ rule.

For an explicit induction invariant fix the input model and proof.
If the label universe has size $h$ after reindexing and the largest
arity at any local sequent or subformula is $k$, there are at most
$|W|\max(1,h)^k$ world-assignment rows per occurrence. Both $h$ and
$k$ come from finite scans, including extra slots under binders.
Compute truth and persistence in bottom-up syntax order, and soundness
in bottom-up proof order. At each stage take the bounded conjunction of
the assertions for all earlier rows. This is a quantifier-free induction
in the language of PR functions, even though the exposition speaks of
every assignment at once.
\end{proof}

\begin{corollary}[Consistency and explicit separation]\label{cor:ded-cons}
Neither calculus has a proof of $\bot$ from the empty context.
There is no $I$ proof of excluded middle $H\vee\neg H$ for a fresh
nullary predicate $H$, whereas $C$ has a proof.
\end{corollary}
\begin{proof}
Take the one-element one-world model with constant functions and empty
relation tables for consistency. For separation, take two worlds $r<s$
with a singleton domain and force $H$ only at $s$. At $r$ neither $H$
nor $\neg H$ holds. Finite soundness excludes an $I$ certificate.
In $I$, if $\neg(H\vee\neg H)$, then $H$ gives a contradiction by the
left injection, yielding $\neg H$. The right injection now gives a
contradiction too. Discharge to obtain $\neg\neg(H\vee\neg H)$ and
apply the extra rule of $C$.
\end{proof}

\subsection{Which quantifier identities have intuitionistic proofs?}

Write $P(x),Q(x)$ for formulas with an indicated slot and $H$ for a
formula independent of that slot. Other parameters may occur throughout.
The following equivalences have uniform $I$ proofs.
\begin{align*}
 \neg\exists x\,P(x)&\ \leftrightarrow\ \forall x\,\neg P(x),\tag{1}\\
 \forall x\forall y\,P(x,y)&\ \leftrightarrow\ \forall y\forall x\,P(x,y),\tag{2}\\
 \exists x\exists y\,P(x,y)&\ \leftrightarrow\ \exists y\exists x\,P(x,y),\tag{3}\\
 \forall x\,(P(x)\wedge Q(x))&\ \leftrightarrow\
      (\forall x\,P(x))\wedge(\forall x\,Q(x)),\tag{4}\\
 \exists x\,(P(x)\vee Q(x))&\ \leftrightarrow\
      (\exists x\,P(x))\vee(\exists x\,Q(x)),\tag{5}\\
 \exists x\,(H\wedge Q(x))&\ \leftrightarrow\ H\wedge\exists x\,Q(x),\tag{6}\\
 (H\to\forall x\,P(x))&\ \leftrightarrow\ \forall x\,(H\to P(x)),\tag{7}\\
 ((\exists x\,P(x))\to H)&\ \leftrightarrow\ \forall x\,(P(x)\to H).\tag{8}
\end{align*}
For (1), a putative $P(u)$ contradicts $\neg\exists x\,P(x)$ by
existential introduction. Generalize $u$. Conversely a witness for
$\exists x\,P(x)$ contradicts its particular negation. For (2), introduce
the target's two fresh parameters, instantiate the input in the opposite
order and generalize in the target order. For (3), open the two input
witnesses in order and introduce the same two in the reversed nesting.
Both operations reverse themselves on exchanging the roles of $x,y$.
For (4), specialize at a fresh parameter and project either component,
generalizing the two results separately. For the converse specialize
both inputs at the same parameter, pair them and generalize. For (5),
open the witness and split its disjunction, introducing the matching
existential and disjunct in each case. In reverse, split the outer
disjunction, open its witness and rebuild the inner disjunction
with that witness. For (6), retain $H$ and introduce $\exists x\,Q(x)$
from the witness, or, in reverse, carry $H$ into the witness branch
and pair it with $Q$. For (7), take a fresh $u$ and assume $H$.
Use the implication and specialize its universal result. In reverse,
assume $H$, specialize at fresh $u$, apply and generalize.
For (8), a given $P(u)$ supplies an existential witness to the left
implication. In reverse, open a given existential at $u$ and apply the
specialized right-hand universal there. Independence of $H$ makes the
parameter conditions in (6)--(8) hold. Each argument is a fixed proof
template, using shifts and openings on input codes.

There are also $I$ proofs of
$\exists x\,\neg P(x)\to\neg\forall x\,P(x)$ and
$H\vee\forall x\,P(x)\to\forall x\,(H\vee P(x))$.
For the first, specialize the supposed universal at the negative
witness. For the second split the disjunction. At each fresh individual
use the left injection if given $H$, and otherwise specialize the
universal and use the right injection.

\begin{proposition}[Two converses separated]\label{prop:quant-krip}
The converses of these last two implications have uniform $C$ proofs.
For a fresh unary $P$ and nullary $H$, neither converse has an $I$ proof.
\end{proposition}
\begin{proof}
In $C$, assume $\neg\forall x\,P(x)$ and, toward a contradiction,
$\neg\exists x\,\neg P(x)$. At any fresh $u$, $\neg P(u)$ gives a
forbidden witness, so $\neg\neg P(u)$. Eliminate double negation and
generalize, contradicting the first assumption. Discharge the second
assumption and eliminate double negation once more. For the other
converse assume $\forall x\,(H\vee P(x))$ and split on $H\vee\neg H$.
In the first branch inject $H$. In the second, arbitrary $u$ gives
$H\vee P(u)$. The left alternative contradicts $\neg H$, and the right
gives $P(u)$. Generalize and use the right injection.

For the first countermodel take a root $r$ with two incomparable
terminal successors $s,t$. All domains are $\{0,1\}$. Force no
$P$ at $r$, just $P(0)$ at $s$ and just $P(1)$ at $t$. The universal
fails at all three worlds, so its negation holds at $r$. Each potential
negative witness is defeated at one successor, so the existential
fails there at the root. For the second take $r<s$, with
$D_r=\{0\}$ and $D_s=\{0,1\}$. Force $P(0)$ throughout, never
$P(1)$, and force $H$ just at $s$. The premise holds at $r$: its
local choice is $P(0)$, and either individual at $s$ satisfies $H$.
At the root $H$ fails and $\forall x\,P(x)$ fails by the later
individual $1$. Apply Theorem~\ref{thm:krip-check} in both cases.
\end{proof}

\paragraph{Lean implementation path for the proof checker.}
Define raw inductives \texttt{Term} (variable, function on a list)
and \texttt{Form} with the constructors above, and Boolean functions
\texttt{termOK} and \texttt{formOK} taking the arity and signature.
One nondependent implementation of substitution carries the binder
depth explicitly and uses the depth algorithm of
Theorem~\ref{thm:subst}. Define \texttt{Sequent} with natural-number
arity, assumption list and conclusion, and a proof node with its sequent,
tag, annotation and ordered list of children. For termination one can
implement tree traversals over postorder arrays: each entry refers only
to earlier entries and the structural test verifies the child sorts and
arities. For a context-dependent calculation, first unroll occurrences
(a subtree repeated by two parents gets two addresses), propagate
arities or substitutions downward, then fold upwards. If the array length
is $N$, an unrolling has depth at most $N$ and branching at most the
maximum declared child-list length $b$. The expression $1+(b+1)+\cdots+(b+1)^N$
bounds its size. Ill-formed inputs return failure before traversal.
Define \texttt{localOK} by cases on the tags listed above, and
\texttt{check} by conjunction with the stored child bits. Prove
\texttt{lift\_comp} and \texttt{open\_subst} from the composition
law. Prove the node equalities for renaming and grafting by cases
on the rule tag, using these two lemmas at quantifiers and equality
elimination. Implement forcing using finite tables, with two extra
loops at a universal node and a bounded witness lookup at an existential
node. Prove the persistence and substitution tables before the
per-rule invariant in Theorem~\ref{thm:krip-check}. For an internal
iPRA implementation, replace inductives and arrays by their natural
codes, and generate the same bounded-fold identities. Lean induction on
the conjunction of all earlier Boolean checks then translates to
quantifier-free induction with the model, formula, proof and other
input codes as parameters.

The companion \texttt{NegativeTranslation.lean} implements the typed
rule trees and proves \texttt{FOLNegative.finite\_sound} and
\texttt{FOLNegative.consistency}. Its argument-list datatype permits
function symbols of arbitrary finite arity. Scope and the signature's
arity checks belong to the external syntax encoding. The Lean model
uses $\operatorname{Fin}(n+1)$, arbitrary total operations and decidable
relations. It also proves soundness for a presented nonempty carrier
when the semantic double-negation step is available for each enabled
classical rule. The finite expanding-domain Kripke argument above uses
the same quantifier rules with an additional world parameter.
